\clearpage
\section{Conclusion}
\label{sec:conclusion}

Read together, the five works tell a simple story. Automated QC did not
progress by discovering one perfect outlier formula. It progressed by
building better references for normal behaviour and by combining different
kinds of evidence.

Three conclusions matter most to me.

First, a score, a reference and a threshold should be treated as separate
design choices. A familiar score is not automatically valid when it is moved
from an independent batch to a correlated sensor window.

Second, thresholds need provenance. A number such as 3.5, 30 or a physical
step limit is only meaningful together with the data, sampling rate, window
and trade-off for which it was chosen. The threshold should be documented
and re-evaluated when those conditions change.

Third, combinations beat single detectors. On the synthetic benchmark, a
simple battery of step, variance-persistence and spatial checks detects every
spike, the entire level shift and both frozen-sensor faults, while retaining
partial sensitivity to drift. Each check is transparent enough to explain to
an operator, and together they cover one another's blind spots.

The final lesson is therefore not to automate doubt away. A good QC system
makes doubt explicit: it records which rule objected, preserves the original
measurement, and gives the reviewer enough context to decide what the flag
means. That is the common architecture connecting the five readings and the
principle I would carry back into my own work.

% =====================================================================
% Reproducibility note
% =====================================================================
